\chapter{Interpreting a register value}
\MARSChapterHeading{Interpreting a register value}
\label{revision:lab-integer-representation}

This lab accompanies the main book's chapter on registers and integer
representation. Use its complete \texttt{byte-inspector.asm} program and the
separate \texttt{word-boundary.asm} experiment. Copies are supplied with the book
in \texttt{revision/labs/files}. Work on your own copies so that deliberate
faults do not replace the examples you use for comparison.

The guided byte-inspector lesson presents four predictions, checks five runs
of the complete program, and records three ordered Debug stops. Import the
course's student ZIP through
\textbf{Tools~\(\to\) Import MARS Lab}. The lesson opens in Learn. Answer its
questions there, then use ordinary Run on the imported
\texttt{byte-inspector.asm} when an activity asks for execution. Keep the normal
assembly defaults, including pseudoinstruction support. Enter the requested
integer in the Run console and let the program finish.

The lesson checks the integer consumed by the program and its console output.
A matching pair of printed numbers alone is insufficient. For example, 150 and
406 have the same low byte, but a run given 406 does not establish that you tested
the requested input 150. Read the feedback before repeating a failed run.

Run and Debug observations belong to the source that executed. An edit,
including a comment, requires those observations to be checked again. Use
\textbf{Tools~\(\to\) Continue Downloaded MARS Lab} to return to the first
activity that needs work. Your prediction answers and revealed hints remain.
The guided Debug activity checks the unchanged program's input, mask and
sign-extension boundaries. Resume that Debug session to follow the reporting
code. Learn then opens two independent starters, described below. The deliberate
faults and word-width experiment remain investigations for you to carry out;
completing the lesson does not replace them.

\section{Make two predictions from one pattern}

% Prevent a lone continuation line without changing the page dimensions.
\begingroup
\widowpenalty=10000
Start with input 150. Write it as \texttt{10010110} in binary and \texttt{0x96}
in hexadecimal. Add the unsigned place values to recover 150. Then assign the
highest bit weight $-128$ to obtain the signed byte value $-106$. Your two answers
should come from the same eight bits. No instruction has run yet, and no bits
have changed merely because you chose another interpretation.
\par
\endgroup

Now predict the complete 64-bit results after the program's two conversion
instructions. The mask should leave \texttt{0x0000000000000096} in
\texttt{\$t0}. Sign extension should leave \texttt{0xFFFFFFFFFFFFFF96} in
\texttt{\$t1}. Include the upper digits in your prediction: both registers have
\texttt{0x96} in their low byte, so inspecting only that byte cannot distinguish
the operations.

\section{Observe the boundary that matters}

Select MIPS64 Release 6, set a breakpoint on the \texttt{andi} below
\texttt{interpret}, and start Debug. Enter 150. At the breakpoint, input has
completed, but the mask has not executed. Inspect \texttt{\$v0}, step over the
mask, and inspect \texttt{\$t0}. Step over \texttt{seb} and inspect
\texttt{\$t1}. Record the source line and register alongside each value.
Learn advances after each matching observation; wait for its next instruction
before stepping again. All three stops must come from one Debug session.

A register or memory edit changes the machine independently of its program.
Backstep also changes the execution history you are inspecting. Both are useful
for investigation, but the guided activity records an unmodified forward
execution. Start a fresh Debug session after either operation when you want to
record these three observations. An incomplete sequence must be repeated in
one session; a completed sequence is retained with the current source.

The next highlighted instruction has not executed. If your prediction appears
wrong, first check that you are comparing the same execution boundary. Then
check the selected display format. A hexadecimal full-width register and a signed
decimal byte are different descriptions; comparing their printed text directly
is insufficient.

Resume execution and check the two output numbers. Explain why both results
survive the intervening string-output requests. Locate each instruction that
copies a result into \texttt{\$a0}. This connects the conversion with the value
actually consumed by integer output.

\section{Choose a test that exposes a fault}

In your working copy, replace the \texttt{seb} with a \texttt{move} using the
same registers. Run input 127. The unchanged output does not establish that the
conversion is correct: both instructions produce 127 for this case. Predict what
will happen for 128, then run it.

Debug the changed program with 128 and inspect both result registers immediately
after the copy. Explain why changing a display preference cannot recover a sign
extension that the program never performed. Restore \texttt{seb}, then separately
change the second integer-output argument to copy \texttt{\$t0}. The signed
conversion should now be correct internally while the second output is wrong.
Use the debugger to identify the first boundary at which those two faulty
programs differ. Restore the original before continuing.

Test inputs 0, 127, 128 and 255. State what each contributes: zero, the largest
positive signed byte, the first negative signed pattern, and the largest unsigned
byte. Explain why testing only values below 128 misses the conversion fault.

\section{Separate masking from validation}

Predict the masked result for 256, then run it as a deliberate experiment outside
the program's stated input range. The low byte is zero. The program has discarded
the ninth bit; it has not reported an invalid input.

Record that distinction explicitly. A precondition describes an input the program
expects. A mask computes a result for the bits it receives. A range check would
need instructions that decide whether to accept the input, which this program
does not contain. Do not use this experiment to conclude that every integer
outside the exercise's range is accepted by the console-input service.

\section{Compare arithmetic widths}

Open \texttt{word-boundary.asm} and break at \texttt{inspect}. Both additions
have completed at this boundary. In hexadecimal, compare the full contents of
\texttt{\$t1} and \texttt{\$t2}; in signed 64-bit decimal, compare their meanings.
Explain the upper eight hexadecimal digits before looking at the low eight,
which are identical.

Replacing \texttt{daddiu} with \texttt{addiu} should make the two results agree.
Predict which result changes and why, then test the replacement in a copy.
The source and immediate are unchanged. The result width changes. Service 1
prints only a signed 32-bit integer, so using that service for both registers
would hide the full-width difference this experiment is designed to expose.

\section{Work without the conversion recipe}

Open \texttt{independent/plus-one.asm} from the student package. Its strings
are supplied, but its body only exits. Write the input, calculation and output
instructions before that exit request, or replace the body with your own complete
program. Keep the original unsigned and signed outputs, then add a third line
with prefix \texttt{Unsigned byte plus one: } and the unsigned value plus one.
For input 255, this third result must be 256.

Keep that solution and work separately in
\texttt{independent/wrapped-byte.asm}. Narrow the sum to its low eight bits
before printing the third line with prefix \texttt{Wrapped byte plus one: }.
For input 255, that result must be zero. Keep the original first two outputs
correct in both programs. Choose your own labels and temporary registers; the
public checks compare the console interface.

Before running, name the register that retains each result and the instruction
that supplies each output argument. If a test fails, inspect the value after its
calculation and again before output. Use those observations to decide whether the
fault lies in arithmetic, narrowing or argument preparation. Explain your final
choice of tests and the property each test distinguishes.

\subsection{Check the result that distinguishes the tasks}

Use the public cases 0, 1, 127, 128, 150, 254 and 255. For 254, the third result
is 255 in both programs; for 255, the ordinary sum is 256 and the narrowed sum
is zero. A check that only exercises the smaller inputs cannot distinguish
these contracts. Explain the retained carry bit in the first program and the
operation that discards it in the second.

After the worked observations, Learn opens each independent starter in turn.
Keep the requested file selected and use \textbf{Check this program}. The check
uses your current editor text, including unsaved changes. Each public case runs
in its own simulator, so checking leaves your ordinary Run or Debug session
available for investigation. Read a failed case's expected and actual output,
then trace that input yourself before changing the program.

To reopen either program directly, use
\textbf{Tools~\(\to\) Choose Downloaded MARS Lab Activity} and select it.
\textbf{Continue Downloaded MARS Lab} returns to the last selected activity.
Both routes share the same checks and history.

Your instructor can control independent work, hints and the number of assessments.
One completed set of cases counts as one assessment, whether it passes or fails.
Assembly errors, cancelled checks and results made outdated by an edit do not
consume an attempt. Editing a completed program requires a new check; the lesson
retains the attempt and hint history. A pass establishes the declared cases,
while your explanation must account for the whole input range.

The \texttt{independent} folder also contains a public input/output contract and a
Kotlin check file for each task. These optional checks use JDK 25 and the
course's MARS command-line JAR. Save your source, open a terminal in that folder,
and replace the example JAR path with its actual location:

\begin{verbatim}
java -jar "/path/to/mars-cli.jar" grade \
  plus-one.kts plus-one.asm
java -jar "/path/to/mars-cli.jar" grade \
  wrapped-byte.kts wrapped-byte.asm
\end{verbatim}

The supplied strings define the prompt, prefixes and newlines; console input
echo is separate from program output. Each case must terminate normally within
20,000 instructions. The checks accept equivalent instruction sequences and do
not inspect scratch registers or source labels.
